% The Slide Definitions
%document
\documentclass[10pt]{beamer}
%theme
\usetheme{metropolis}
% packages
\usepackage{color}
\usepackage{listings}
\usepackage[ngerman]{babel}
\usepackage[utf8]{inputenc}
\usepackage{multicol}


% color definitions
\definecolor{mygreen}{rgb}{0,0.6,0}
\definecolor{mygray}{rgb}{0.5,0.5,0.5}
\definecolor{mymauve}{rgb}{0.58,0,0.82}

\lstset{
    backgroundcolor=\color{white},
    % choose the background color;
    % you must add \usepackage{color} or \usepackage{xcolor}
    basicstyle=\footnotesize\ttfamily,
    % the size of the fonts that are used for the code
    breakatwhitespace=false,
    % sets if automatic breaks should only happen at whitespace
    breaklines=true,                 % sets automatic line breaking
    captionpos=b,                    % sets the caption-position to bottom
    commentstyle=\color{mygreen},    % comment style
    % deletekeywords={...},
    % if you want to delete keywords from the given language
    extendedchars=true,
    % lets you use non-ASCII characters;
    % for 8-bits encodings only, does not work with UTF-8
    frame=single,                    % adds a frame around the code
    keepspaces=true,
    % keeps spaces in text,
    % useful for keeping indentation of code
    % (possibly needs columns=flexible)
    keywordstyle=\color{blue},       % keyword style
    % morekeywords={*,...},
    % if you want to add more keywords to the set
    numbers=left,
    % where to put the line-numbers; possible values are (none, left, right)
    numbersep=5pt,
    % how far the line-numbers are from the code
    numberstyle=\tiny\color{mygray},
    % the style that is used for the line-numbers
    rulecolor=\color{black},
    % if not set, the frame-color may be changed on line-breaks
    % within not-black text (e.g. comments (green here))
    stepnumber=1,
    % the step between two line-numbers.
    % If it's 1, each line will be numbered
    stringstyle=\color{mymauve},     % string literal style
    tabsize=4,                       % sets default tabsize to 4 spaces
    % show the filename of files included with \lstinputlisting;
    % also try caption instead of title
    language = [Sharp]C,
	showspaces = false,
	showtabs = false,
	showstringspaces = false,
	escapechar = ,
}

\def\ContinueLineNumber{\lstset{firstnumber=last}}
\def\StartLineAt#1{\lstset{firstnumber=#1}}
\let\numberLineAt\StartLineAt



\newcommand{\codeline}[1]{
	\alert{\texttt{#1}}
}


% Author and Course information
% This Document contains the information about this course.

% Authors of the slides
\author{\href{https://github.com/evilham}{@evilham}\\
\footnotesize{Basierend auf Folien von \href{https://github.com/satkowski}{@satkowski} (November 2016)}}

% Name of the Course
\institute{C\texttt{\#} Kurs}

% Fancy Logo
\titlegraphic{\hfill\includegraphics[height=1.25cm]{../templates/fsr_logo_cropped}}


% Presentation title
\title{Weiteres ins C\texttt{\#}}
\date{\today}


\begin{document}

\maketitle

\begin{frame}{Gliederung}
	\setbeamertemplate{section in toc}[sections numbered]
	\tableofcontents
\end{frame}

% ----------------------- Eventhandling ------------------------------
\section{Event Handling}
\subsection{Event Handler}
\begin{frame}{EventHandler}
	\begin{itemize}
		\item Ein vordefinierter Delegate
		\item Wird genutzt um Eventhandling zu betreiben
		\begin{itemize}
			\item Man kann auch einen eigenen nehmen
		\end{itemize}
		\item Gibt \alert{void} zurück
		\item Hat 2 Übergabeparameter:
		\begin{itemize}
			\item \alert{object} (der Sender dieses Events)
			\item \alert{TEventArgs}
		\end{itemize}
	\end{itemize}
	\begin{itemize}
		\item Der Kopf des EventHandlers im Framework
	\end{itemize}
	\lstinputlisting{resources/10_weiteres/event_handler.cs}	
\end{frame}

\begin{frame}{EventArgs}
	\begin{itemize}
		\item Gibt dem Event die Möglichkeit zusätzliche Informationen an Empfänger zu übergeben
		\begin{itemize}
			\item Falls keine Informationen zum Event zur Verfügung steht, kann man \alert{EventArgs.Empty} übergeben
		\end{itemize}
		\item Da der EventHanlder generisch ist, kann man seine eigenen \alert{EventArgs} übergeben
	\end{itemize}
	\lstinputlisting{resources/10_weiteres/event_args.cs}
\end{frame}

\subsection{event}
\begin{frame}{event}
	\begin{itemize}
		\item Wird genutzt um bei einem Publisher ein Event zu deklarieren
		\item Zusammen mit Attributen von einem Delegatetypen genutzt
		\item Kann ebenfalls alle Modifikatoren (wie \alert{virtual} oder \alert{override} nutzen)
		\begin{itemize}
			\item Bei \alert{abstract} muss ein eigener add/remove-Ereignisaccessorblock erstellt werden
		\end{itemize}
	\end{itemize}				
	\lstinputlisting{resources/10_weiteres/event.cs}
\end{frame}

\subsection{Mehrere Events}
\begin{frame}{Mehrere Events}
	\begin{itemize}
		\item Man kann verschiedenenen EventHandlern Methoden übergeben, etwa mit dem add-Ereignisaccessorblock
			\begin{itemize}
				\item \texttt{\alert{+=}}
			\end{itemize}
		\item Entfernt können diese wieder mit dem remove-Ereignisaccessorblock
			\begin{itemize}
				\item \texttt{\alert{-=}}
			\end{itemize}		
		\item Die Methoden werden dann nach einander ausgeführt
	\end{itemize}
\end{frame}

\subsection{Beispiel}
\begin{frame}{Beispiel}
	\lstinputlisting{resources/10_weiteres/event_handling_creation.cs}
\end{frame}

\begin{frame}{Beispiel}
	\lstinputlisting{resources/10_weiteres/event_handling_usage_1.cs}
\end{frame}

\begin{frame}{Beispiel}
	\lstinputlisting{resources/10_weiteres/event_handling_usage_2.cs}
\end{frame}

% ----------------------- Weiteres zu null ------------------------------
\section{Weiteres zu null}
\subsection{Nullable}
\begin{frame}{Nullable}
	\begin{itemize}
		\item Werden genutzt, damit Werttyp-Variablen ebenfalls \alert{null} sein können
		\item Generische Klasse
		\item Besitzt 2 Eigenschaften
		\begin{description}
			\item[HasValue] Gibt an, ob das Objekt einen Wert hat oder nicht
			\item[Value] Gibt den Wert des Objektes zurück\\ Wirft eine InvalidOperationException falls HasValue \alert{false} ist 
		\end{description}
		\item Syntax:
		\begin{itemize}
			\item \texttt{\alert{<Werttyp>}? \alert{<Variablenname>};} oder
			\item \texttt{Nullable<\alert{<Datentyp>}> \alert{<Variablenname>};}
		\end{itemize}
	\end{itemize}
\end{frame}

\begin{frame}{Nullable}
	\lstinputlisting{resources/10_weiteres/nullable.cs}	
\end{frame}

\begin{frame}{Weiteres zu null}
	\begin{itemize}
		\item Falls ein Objekt \alert{null} sein kann sollte man diese vor dem benutzen immer überprüfen
	\end{itemize}
	\begin{itemize}	
		\item \alert{obj != null} kann durch \texttt{\textbf{\alert{?.}}} vereinfacht werden	
		\item Syntax:
		\begin{itemize}
			\item \texttt{obj\textbf{\alert{?.}}Property}
		\end{itemize}
		\item Es wird vor dem Benutzen des Objektes geschaut, ob dieses \alert{null} ist
		\item Falls es nicht \alert{null} ist, wird es benutzt, sonst gibts es \alert{null} zurück
	\end{itemize}
\end{frame}

\begin{frame}{Weiteres zu null}
	\lstinputlisting{resources/10_weiteres/null1.cs}
	\begin{itemize}
		\item Hier für einen \alert{delegate}
	\end{itemize}		
	\lstinputlisting{resources/10_weiteres/null2.cs}	
\end{frame}

% ----------------------- LINQ ------------------------------
\section{LINQ}
\begin{frame}{LINQ}
	\begin{itemize}
		\item Language-Integrated Query
		\item LINQ ermöglicht es Abfragen an Objekte und Datenstrukturen zu stellen
		\item Ähnlich zu z.B. SQL
	\end{itemize}
\end{frame}

\begin{frame}{LINQ}
	\lstinputlisting{resources/10_weiteres/linq.cs}	
\end{frame}

\end{document}
